\documentclass[12pt]{article}
\usepackage[T2A]{fontenc}
\usepackage[utf8]{inputenc}
\usepackage[russian]{babel}
\usepackage{amsmath, amssymb, amsfonts}
\usepackage{geometry}
\usepackage{fancyhdr}
\usepackage{microtype}
\usepackage{titlesec} % Пакет для управления стилем заголовков
\usepackage{enumitem}
\usepackage{tikz}
\usepackage{tikz-3dplot}
\usetikzlibrary{calc, intersections, through, backgrounds, arrows.meta}

% Настройка Section: по центру, жирный, крупный
\pagestyle{fancy}
% Настройка полей
\geometry{a5paper, margin=1cm, top=2cm, headheight=15pt}

% --- ЖЕСТКАЯ КОРРЕКЦИЯ ЦЕНТРИРОВАНИЯ ---
\pagestyle{fancy}
\fancyhf{} % Полная очистка всех полей (убирает невидимые отступы)
\fancyhead[C]{\thepage} % Номер строго по центру
\renewcommand{\headrulewidth}{0pt} % Убираем линию
\fancyfoot{} % Низ пустой

% Убираем внутренние отступы колонтитулов, которые могут давать перекос
\setlength{\headwidth}{\textwidth} 

% Переопределяем системный стиль plain, чтобы он не сбивал настройки
\makeatletter
\let\ps@plain\ps@fancy
\makeatother

\parindent=0pt
\parskip=0.6em
\titleformat{\section}
  {\normalfont\bfseries\centering}{}{0pt}{}

% Настройка Subsection: отступ 1см, жирный, обычный размер
% {команда}[формат]{стиль}{метка}{отступ_от_метки}{перед_заголовком}[после_заголовка]
\titleformat{\subsection}
  {\normalfont\bfseries}{\hspace{1cm}}{1pt}{}
\titlespacing*{\subsection}{1cm}{2ex plus 1ex minus .2ex}{0.5ex plus .2ex}

% Убираем автоматическую нумерацию, чтобы заголовки были как в оригинале
\setcounter{secnumdepth}{0}
\AddEnumerateCounter{\asbuk}{\russian@asbuk@alph}{щ} % section 8.1
\setlist{wide, nosep, font=\rmfamily\normalfont}
\setlist[enumerate, 2]{label=\asbuk*)}
\setlist[enumerate, 3]{label=\asbuk*)}
